% Updated 30 August 2026
\documentclass[10pt]{article}
\usepackage[margin=0.62in]{geometry}
\usepackage[T1]{fontenc}
\usepackage{lmodern,microtype}
\usepackage{amsmath,amssymb,amsthm,mathtools}
\usepackage{enumitem,booktabs,tabularx,array,multicol,xcolor}
\usepackage[bookmarks=false]{hyperref}
\hypersetup{colorlinks=true,linkcolor=blue!45!black,urlcolor=blue!45!black}
\setlength{\parindent}{0pt}
\setlength{\parskip}{3pt}
\setlist{leftmargin=*,itemsep=1pt,topsep=2pt}
\newcommand{\F}{\mathbb F}
\newcommand{\Prb}{\mathbb P}
\newcommand{\E}{\mathbb E}
\newcommand{\IP}{\mathsf{IP}}
\newcommand{\MIP}{\mathsf{MIP}}
\newcommand{\PSPACE}{\mathsf{PSPACE}}
\newcommand{\NEXP}{\mathsf{NEXP}}
\newcommand{\PH}{\mathsf{PH}}
\newcommand{\AM}{\mathsf{AM}}
\newcommand{\NP}{\mathsf{NP}}
\newcommand{\coNP}{\mathsf{coNP}}
\newtheorem*{theorem}{Theorem}
\newtheorem*{lemma}{Lemma}
\newenvironment{keybox}{\begin{quote}\small\itshape}{\end{quote}}

\begin{document}
\begin{center}
{\LARGE\bfseries Interactive Proof Systems\par}
\vspace{1mm}
{Hadi Zamani \qquad Examiner: Michal Kouck\'y\par}
{\small 30 August 2026}
\end{center}

\section*{Assigned articles}
\begin{itemize}
\item \emph{Algebraic Methods for Interactive Proof Systems} (LFKN)
\item \emph{IP = PSPACE} (Shamir), with \emph{IP = PSPACE: Simplified Proof} (Shen)
\item \emph{Non-Deterministic Exponential Time Has Two-Prover Interactive Protocols} (BFL)
\item \emph{Proofs that Yield Nothing but Their Validity or All Languages in NP
Have Zero-Knowledge Proof Systems} (GMW)
\end{itemize}

\subsection*{Article map: main theorems and techniques}
\begingroup
\scriptsize
\setlength{\tabcolsep}{3.5pt}
\renewcommand{\arraystretch}{1.08}
\begin{center}
\begin{tabularx}{\textwidth}{@{}>{\raggedright\arraybackslash}p{0.18\textwidth}>{\raggedright\arraybackslash}X>{\raggedright\arraybackslash}X@{}}
\toprule
\textbf{Assigned article / note} & \textbf{Main theorem or result} &
\textbf{Characteristic techniques}\\
\midrule
LFKN (1992) & Algebraic protocols verify permanent/$\#\mathsf P$ claims and imply
$\PH\subseteq\IP$. & Arithmetization; low-degree self-reducibility; sum-check and
many-values-to-one random reduction; Schwartz--Zippel soundness.\\
Shamir (1992) & $\IP=\PSPACE$: quantified Boolean formulas have interactive
proofs, while protocol game trees are evaluable in polynomial space. & TQBF
arithmetization; simple-QBF preprocessing and dummy variables for degree control;
fresh random low-degree checks; depth-first exact game-tree evaluation.\\
Shen (1992), companion note & Gives a simplified proof of
$\PSPACE\subseteq\IP$. & Explicit multilinearization operators $R_i$; local
interpolation identities; controlled individual degree instead of simple-QBF syntax.\\
BFL (1991) & $\MIP=\NEXP$. & Arithmetize an exponential tableau; low-degree
extension and multilinearity tests; cross-check noncommunicating provers; LFKN-style
global constraint verification.\\
GMW (1991) & Assuming suitable commitments, every language in $\NP$ has a
zero-knowledge proof; 3-colouring is the canonical protocol. & Hiding/binding
commitments; random colour permutation; random-edge challenge; rewinding simulation;
sequential repetition.\\
\bottomrule
\end{tabularx}
\end{center}
\endgroup

\section*{Theorem inventory}
\begingroup
\small

\subsection*{Lund--Fortnow--Karloff--Nisan (LFKN)}
\begin{itemize}
\item \textbf{Theorem 1.} Every language in $\mathsf P^{\#\mathsf P}$ has an
interactive proof system.
\item \textbf{Corollary 2.} By Toda's theorem,
$\PH\subseteq\IP$; in particular, $\coNP\subseteq\IP$.
\item \textbf{Lemma 3.} An interactive protocol for verifying the permanent of
a $0$--$1$ matrix suffices to give interactive proofs for every language in
$\mathsf P^{\#\mathsf P}$.
\item \textbf{Lemma 4.} If the claimed polynomial along a random matrix line is
not the true permanent polynomial, a uniformly random field point catches the
discrepancy except with probability at most $r/|\F|$.
\item \textbf{Corollary 5.} If cryptographic one-way functions exist, every
language in $\PH$ has a zero-knowledge interactive proof.
\item \textbf{Corollary 6.} If every language in $\IP$ had a bounded-round
interactive proof, then the polynomial hierarchy would collapse.
\item \textbf{Corollary 7.} The permanent has a probabilistic polynomial-time
program checker that either certifies the program's answer on the given matrix
or finds evidence that the program fails somewhere.
\item \textbf{Corollary 8.} If $\#\mathsf P$ has polynomial-size circuits, then
$\mathsf P^{\#\mathsf P}\subseteq\mathsf{MA}$.
\end{itemize}

\subsection*{Shamir: \emph{IP = PSPACE}}
\begin{itemize}
\item \textbf{Theorem 1.} Every QBF of size $n$ can be transformed into an
equivalent simple QBF of size polynomial in $n$.
\item \textbf{Theorem 2.} A closed QBF is true if and only if the value of its
arithmetization is nonzero.
\item \textbf{Theorem 3.} The integer value of the arithmetization of a closed
size-$n$ QBF is at most doubly exponential, $2^{2^{O(n)}}$.
\item \textbf{Theorem 4.} If the arithmetized value is nonzero, there is a prime
of polynomial bit length modulo which it remains nonzero.
\item \textbf{Theorem 5.} For a simple QBF, every univariate polynomial queried
in the protocol has degree at most linear in the formula size; dummy variables
can reduce the relevant degree to a constant.
\item \textbf{Theorem 6.} The QBF protocol has perfect completeness, while on a
false QBF every cheating prover is accepted only with negligible probability.
\item \textbf{Theorem 7.} Every $\PSPACE$ language is accepted even by Shamir's
weak verifier model.
\item \textbf{Conclusion.} Theorems 1--6 prove
$\PSPACE\subseteq\IP$; together with the previously known
$\IP\subseteq\PSPACE$, they give $\IP=\PSPACE$.
\end{itemize}

\subsection*{Shen: simplified proof}
\begin{itemize}
\item \textbf{Main theorem.} $\PSPACE\subseteq\IP$.  Shen replaces Shamir's
simple-QBF preprocessing by explicit multilinearization operators $R_i$ and
checks the resulting local polynomial identities at fresh random points.
\end{itemize}

\subsection*{Babai--Fortnow--Lund (BFL)}
\begin{itemize}
\item \textbf{Theorem 1.1.} $\MIP=\NEXP$; two-prover interactive proofs have
exactly nondeterministic exponential-time power.
\item \textbf{Theorem 2.3 (Fortnow--Rompel--Sipser).} Probabilistic oracle
machines and multiprover interactive protocols recognize the same languages.
\item \textbf{Theorem 2.4 (Ben-Or--Goldwasser--Kilian--Wigderson).} Any
multiprover protocol can be simulated with two noncommunicating provers.
\item \textbf{Theorem 4.6.} The multilinearity tester accepts a genuinely
multilinear, bounded integer-valued oracle with certainty and rejects an oracle
far from multilinear with high probability.
\item \textbf{Corollary 4.10.} Every language in $\mathsf{EXP}$ has a
multiprover interactive proof with $\mathsf{EXP}$-complexity provers.
\item \textbf{Theorem 5.13.} If only an $\varepsilon$ fraction of the axis-parallel
lines are wrong, the oracle is globally
$\varepsilon_0$-close to a multilinear function, with
$\varepsilon_0=3n^2(\varepsilon+\delta+1/N)$ under the stated parameters.
\item \textbf{Corollary 6.3.} Every $\mathsf{EXP}$-complete language has a
function-restricted interactive proof and therefore a program checker.
\item \textbf{Theorem 6.4.} Every $\PSPACE$-complete and
$\mathsf{EXP}$-complete language has a self-testing/self-correcting pair.
\item \textbf{Corollary 6.7.} There are $\PSPACE$-complete and
$\mathsf{EXP}$-complete functions with self-testing/self-correcting pairs in
the Blum--Luby--Rubinfeld model.
\item \textbf{Theorem 6.8 (Meyer--Karp--Lipton).} For
$\mathcal C\in\{\mathsf{EXP},\PSPACE,\mathsf P^{\#\mathsf P}\}$,
$\mathcal C\subseteq\mathsf P/\mathsf{poly}$ implies
$\mathcal C=\Sigma_2^{\mathsf P}$.
\item \textbf{Corollary 6.9.} If a language has a multiprover proof with
provers of complexity $\mathcal C$, and $\mathcal C$ has polynomial-size
circuits, then the language lies in $\mathsf{MA}$.
\item \textbf{Corollary 6.10.} If all of $\mathsf{EXP}$ has polynomial-size
circuits, then
$\mathsf{EXP}=\Sigma_2^{\mathsf P}=\Pi_2^{\mathsf P}=\mathsf{MA}$; analogously
for $\PSPACE$ and $\mathsf P^{\#\mathsf P}$.
\end{itemize}

\subsection*{Goldreich--Micali--Wigderson (GMW)}
\begin{itemize}
\item \textbf{Theorem 1.} Protocol 1 is a two-message interactive proof for
Graph Nonisomorphism.
\item \textbf{Corollary 1.} Relative to a random oracle, Graph Isomorphism lies
in $\NP\cap\coNP$; nonuniformly, Graph Nonisomorphism lies in
$\NP/\mathsf{poly}$.
\item \textbf{Corollary 2.} If Graph Isomorphism were NP-complete, the
polynomial hierarchy would collapse to its second level.
\item \textbf{Theorem 2.} Protocol 2 is a perfect zero-knowledge interactive
proof for Graph Isomorphism.
\item \textbf{Theorem 3.} The verifier-tested modification, Protocol 3, is a
perfect zero-knowledge interactive proof for Graph Nonisomorphism.
\item \textbf{Proposition 4.} Assuming a nonuniformly secure encryption
function, Protocol 4 is a computational zero-knowledge proof for Graph
$3$-Colourability.
\item \textbf{Lemma 4.0.} Polynomially many encryptions of any two message
sequences are computationally indistinguishable to polynomial-size circuit
families; this is the hybrid lemma used by the simulator.
\item \textbf{Theorem 5.} If nonuniformly secure encryption exists, every
language in $\NP$ has a zero-knowledge interactive proof.
\item \textbf{Theorem 6.} The preceding theorem can use a probabilistic
polynomial-time prover supplied with an NP witness as auxiliary input.
\item \textbf{Simulation claims 2.1--2.2, 3.1--3.2, and 4.1--4.2.} In each
protocol the constructed simulator runs in expected polynomial time and its
output is respectively identical to, or computationally indistinguishable from,
the cheating verifier's real view.
\end{itemize}

\endgroup

\section{Definitions, quantifiers, and amplification}

An interactive proof for $L$ is a probabilistic polynomial-time verifier $V$
interacting for polynomially many rounds with an unbounded prover.  With the
standard constants,
\[
 x\in L\Longrightarrow \exists P\quad
   \Prb[(P,V)(x)\text{ accepts}]\geq\tfrac23,
\]
\[
 x\notin L\Longrightarrow \forall P^*\quad
   \Prb[(P^*,V)(x)\text{ accepts}]\leq\tfrac13.
\]
The order $\exists P$ in completeness and $\forall P^*$ in soundness is
essential.  A proof has soundness against computationally unbounded provers;
an \emph{argument} need only resist polynomial-time cheating provers.

Sequential repetition followed by majority reduces both errors exponentially;
fresh verifier randomness is required.  Public-coin protocols expose each random
challenge.  Private coins do not increase the power of polynomial-round
interactive proofs (Goldwasser--Sipser), although the conversion is not the
content of the four main articles.  For a fixed number $k\geq2$ of public-coin
Arthur--Merlin rounds, the standard round-collapse gives $\AM(k)=\AM(2)$.
This is an internal collapse of that protocol hierarchy, not a collapse of the
polynomial hierarchy.

\section{The algebraic engine of LFKN}

\subsection{Arithmetization and polynomial identity testing}

Over a field $\F$, arithmetize Boolean operations by
\[
 \neg x\mapsto1-x,\qquad x\wedge y\mapsto xy,
 \qquad x\vee y\mapsto x+y-xy.
\]
These agree with Boolean logic on $\{0,1\}$; outside that cube they are simply
polynomials.  For $f:\{0,1\}^m\to\F$, its unique multilinear extension is
\[
 \widetilde f(z)=\sum_{a\in\{0,1\}^m} f(a)
 \prod_{i=1}^m\bigl(a_i z_i+(1-a_i)(1-z_i)\bigr).
\]

\begin{lemma}[Schwartz--Zippel]
If $0\neq P\in\F[x_1,\ldots,x_m]$ has total degree at most $d$ and
$r$ is uniform in $S^m\subseteq\F^m$, then
\[
 \Prb[P(r)=0]\leq d/|S|.
\]
\end{lemma}
Proof by induction on $m$: expand in $x_m$; with probability at most
$(d-k)/|S|$ its leading coefficient vanishes, and otherwise the resulting
univariate polynomial has at most $k$ roots.  The lemma is the source of
soundness: two distinct low-degree polynomials rarely agree at a fresh point.

\subsection{Sum-check: statement, protocol, and soundness}

The prover claims
\[
 H=\sum_{b\in\{0,1\}^m}g(b),
\]
where the individual degree of $g$ is at most $d$ and the verifier can evaluate
$g$ at one field point efficiently.  The protocol is:
\begin{enumerate}
\item The prover sends
$q_1(X)=\sum_{b_2,\ldots,b_m}g(X,b_2,\ldots,b_m)$.
Check $\deg q_1\leq d$ and $q_1(0)+q_1(1)=H$; choose uniform $r_1\in\F$.
\item In round $i$, receive $q_i$, check $\deg q_i\leq d$ and
\[
 q_i(0)+q_i(1)=q_{i-1}(r_{i-1}),
\]
then choose fresh uniform $r_i\in\F$.
\item Check $q_m(r_m)=g(r_1,\ldots,r_m)$ directly.
\end{enumerate}
An honest prover passes with probability one.  At the first dishonest round,
the sent univariate and the correct one are distinct.  They agree at the fresh
$r_i$ with probability at most $d/|\F|$.  A union bound gives soundness error
at most $md/|\F|$.  The verifier reads $O(md)$ field elements rather than
$2^m$ summands.

\subsection{LFKN and the permanent}

For $A\in\F^{n\times n}$,
\[
 \operatorname{per}(A)=\sum_{\pi\in S_n}\prod_{i=1}^n A_{i,\pi(i)}
 =\sum_{j=1}^n A_{1j}\operatorname{per}(A_{1j}),
\]
where $A_{1j}$ deletes row $1$ and column $j$.  Downward self-reducibility is
easy, but expanding the entire tree is exponential.  LFKN places related
subinstances on a low-degree algebraic curve.  A claimed polynomial containing
their permanent values is checked at the required points; the verifier then
chooses a fresh random point and reduces many old claims to one new claim on a
smaller instance.  Low degree plus Schwartz--Zippel prevents the prover from
changing the polynomial after seeing the point.

Because the permanent is $\#\mathsf P$-complete and Toda's theorem gives
$\PH\subseteq\mathsf P^{\#\mathsf P}$, the LFKN protocol yields
\[
 \PH\subseteq\IP.
\]
The conceptual payload is more reusable than this particular application:
arithmetize, maintain a low-degree claim, reduce many values to one random
value, and finish at an explicitly computable base case.

\section{$\IP=\PSPACE$: both inclusions in full outline}

\subsection{$\IP\subseteq\PSPACE$: exact game-tree evaluation}

Normalize the protocol so that nodes alternate between a prover message and a
block of verifier coin tosses.  For a transcript prefix $h$, define its value by
\[
 \operatorname{Val}(h)=
 \begin{cases}
  0\text{ or }1, &h\text{ is a leaf},\\
  \max_m\operatorname{Val}(hm), &h\text{ is a prover node},\\
  2^{-r}\sum_{z\in\{0,1\}^r}\operatorname{Val}(hz),
      &h\text{ is an }r\text{-coin verifier node}.
 \end{cases}
\]
There are exponentially many children, but every message has polynomial length
and the depth is polynomial.  A deterministic machine enumerates the children
with a polynomial-bit counter and evaluates them depth first, reusing memory.
The recursion stack stores only polynomially many polynomial-size transcripts.
Values are exact rationals with denominator at most $2^R$, where the total
number $R$ of verifier coins is polynomial, so their numerators and denominators
also use polynomially many bits.  Comparing the root value with the completeness
threshold therefore uses polynomial space.  This proves $\IP\subseteq\PSPACE$;
finiteness of the tree alone is not a sufficient space analysis.

\subsection{Why TQBF is $\PSPACE$-complete}

Membership is by depth-first recursive evaluation of the quantified formula,
storing one assignment and one recursion path.  For hardness, let $M$ use
$s(n)$ space.  A configuration has $O(s(n))$ bits, and if $M$ accepts then it
does so within $2^{O(s(n))}$ steps after standard normalization.

Let $R_i(C,D)$ mean that $D$ is reachable from $C$ within $2^i$ steps.
The na\"ive recurrence duplicates $R_i$; the following quantified selector keeps
one copy:
\[
\begin{split}
R_{i+1}(C,D):={}&\exists M\ \forall b\ \exists U,V\;[
 (b=0\Rightarrow(U=C\wedge V=M))\\[-1mm]
&\hspace{36mm}\wedge
 (b=1\Rightarrow(U=M\wedge V=D))
 \wedge R_i(U,V)].
\end{split}
\]
Here $M,U,V$ are configuration-sized bit vectors, and $R_0$ checks equality or
one legal transition.  The universal bit chooses which half of the path must be
verified; because $U,V$ are quantified after $b$, they may depend on that choice.
Crucially, the displayed formula contains only one syntactic copy of $R_i(U,V)$, so
\[
 |R_{i+1}|=|R_i|+O(s(n)),
\]
instead of the doubling caused by writing both halves explicitly.  With
$i=O(s(n))$, the QBF has polynomial size and is true exactly when an accepting
configuration is reachable.

\subsection{$\PSPACE\subseteq\IP$: quantified arithmetization}

It suffices to give an interactive proof for
\[
 \Phi=Q_1x_1\cdots Q_nx_n\,\varphi(x_1,\ldots,x_n).
\]
Let $P$ be the arithmetization of $\varphi$.  Quantifiers become algebraic
operators
\[
 A_iP=P|_{x_i=0}\,P|_{x_i=1},\qquad
 E_iP=P|_{x_i=0}+P|_{x_i=1}-P|_{x_i=0}P|_{x_i=1}.
\]
On Boolean inputs, $A_i$ represents $\forall x_i$ and $E_i$ represents
$\exists x_i$.  Thus $\Phi$ is true iff the resulting scalar is $1$.

The obstacle is degree: products introduced by $A_i,E_i$ can double degrees in
variables not yet eliminated.  Over the Boolean cube one may replace a polynomial
by its degree-one representative in variable $x_i$:
\[
 R_iP=P|_{x_i=0}+x_i(P|_{x_i=1}-P|_{x_i=0})
      \equiv P\pmod{x_i^2-x_i}.
\]
Insert the required $R_i$ operators so that every intermediate individual degree
is polynomially bounded.  The truth of $\Phi$ is now an explicitly ordered,
polynomial-length sequence of local operators $A_i,E_i,R_i$ applied to $P$.

The verifier starts from the claim that this operator expression equals $1$.
At an $A_i$ or $E_i$ step, the prover sends the relevant low-degree univariate
$p(X)$.  The verifier checks the claimed degree and the endpoint identity
\[
 v=p(0)p(1)\quad(A_i),\qquad
 v=p(0)+p(1)-p(0)p(1)\quad(E_i).
\]
At an $R_i$ step it checks the analogous interpolation identity
$v=p(0)+a(p(1)-p(0))$ at the current field coordinate $a$.
It then chooses a fresh uniform field element $r$, replaces the current claim by
$p(r)$, and proceeds to the shorter operator expression.  At the end it evaluates
the explicit formula polynomial $P$ itself at the accumulated field point.

The invariant is that the verifier holds a claimed value of the remaining operator
expression after all earlier random coordinates have been fixed.  If this claim is
false, any polynomial satisfying the endpoint identity but supporting the false value
must differ from the honest low-degree polynomial.  Their difference is nonzero, so a
fresh random evaluation preserves the lie with probability at most $d/|\F|$.  This is
why the prover must send its polynomial before seeing the next challenge.

If all claims are honest, every identity holds.  At the first false univariate
claim, Schwartz--Zippel lets it survive the fresh challenge with probability at
most $d/|\F|$.  If there are $T=\operatorname{poly}(|\Phi|)$ checks and degree
$d=\operatorname{poly}(|\Phi|)$, choose a polynomial-size field with
$|\F|>3Td$; the total error is below $1/3$.  Field elements and arithmetic have
polynomial bit complexity.  Hence $\PSPACE\subseteq\IP$.

\subsection{Shamir versus Shen}

\begin{center}\small
\begin{tabularx}{\textwidth}{@{}lXX@{}}
\toprule
& \textbf{Degree control} & \textbf{Proof organization}\\
\midrule
Shamir & preprocesses to a \emph{simple QBF}, using dummy variables and
constraints to keep the arithmetized degrees polynomial & syntactic preparation
first, then the low-degree protocol\\
Shen & uses explicit degree-reduction operators $R_i$ instead of simple QBFs &
operator sequence and local interpolation checks make the bookkeeping shorter\\
\bottomrule
\end{tabularx}
\end{center}
Both are algebraic, nonrelativizing proofs of the same containment.

\paragraph{The degree-reduction calculation to reproduce orally.}
For a polynomial $q$ and a coordinate $x_i$, Shen uses
\[
 (R_iq)(x_1,\ldots,x_n)
 =(1-x_i)q(x_1,\ldots,0,\ldots,x_n)
   +x_iq(x_1,\ldots,1,\ldots,x_n).
\]
Thus $R_iq$ agrees with $q$ whenever $x_i\in\{0,1\}$, but has individual
degree at most one in $x_i$.  After an arithmetized quantifier step raises a
variable's degree, inserting $R_i$ restores the Boolean values while preventing
degree growth from accumulating.  At a field point whose $i$th coordinate is
$a$, the verifier checks the local identity
\[
 (R_iq)(\ldots,a,\ldots)
 =q(\ldots,0,\ldots)+a\bigl(q(\ldots,1,\ldots)-q(\ldots,0,\ldots)\bigr).
\]
This is the concrete reason Shen can work directly with the original QBF rather
than first converting it into Shamir's simple-QBF normal form.

\subsection{Relativization and nearby class questions}

For every oracle $O$, the game-tree argument still gives
$\IP^O\subseteq\PSPACE^O$.  Equality does \emph{not} hold for every oracle:
there is an oracle $B$ for which
$\coNP^B\nsubseteq\IP^B$, while always $\coNP^B\subseteq\PSPACE^B$; hence
$\IP^B\neq\PSPACE^B$.  In fact, random oracles separate the two with probability
one.  Equality can hold relative to an oracle: if $A$ is a $\PSPACE$-complete
language, then
\[
 \PSPACE=\mathsf P^A\subseteq\IP^A\subseteq\PSPACE^A=\PSPACE.
\]
Thus the theorem has contradictory relativizations, a standard sign that its
arithmetization is genuinely nonrelativizing.

\section{BFL: two provers and $\MIP=\NEXP$}

In classical $\MIP$, two or more all-powerful provers agree on a strategy before
the protocol but cannot communicate after receiving their questions.  Their
answers must therefore be marginals of fixed global strategy functions.  This
noncommunication is the extra resource: with one prover, every claimed value can
be adapted to the whole transcript; a second prover can be asked a correlated
question that tests consistency with one global object.

\subsection{$\MIP\subseteq\NEXP$}

Fix the input.  A deterministic strategy for each prover is a table mapping its
polynomial-length local transcripts to polynomial-length answers.  There are at
most $2^{\operatorname{poly}(n)}$ entries.  A nondeterministic exponential-time
machine guesses all tables, enumerates all polynomial-length verifier random
strings, simulates the protocol, and counts acceptance exactly.  If randomized
provers helped, fixing their private coins would leave some deterministic strategy
with at least the same acceptance probability.  Hence $\MIP\subseteq\NEXP$.

\subsection{$\NEXP\subseteq\MIP$: proof architecture}

An $\NEXP$ computation has an exponentially long tableau or witness, while every
location still has a polynomial-length address.  The lower-bound proof has four
separate stages.
\begin{enumerate}
\item Encode the accepting computation by a Boolean function on
      $m=\operatorname{poly}(n)$ address bits and extend it to the unique
      multilinear polynomial over a sufficiently large finite field.
\item Arithmetize the local transition constraints.  A false tableau makes the
      resulting constraint polynomial nonzero; an LFKN-style reduction turns the
      claim that its global sum is zero into claims about values at random points.
\item Replace trusted oracle access by two noncommunicating provers.  Correlated
      questions compare their answers, while the multilinearity test checks that
      the alleged value table behaves like one multilinear function rather than
      unrelated answers chosen separately at each query.
\item For soundness, first condition on passing the multilinearity and consistency
      tests.  The answers are then close to evaluations of a single multilinear
      polynomial.  If no accepting tableau exists, the arithmetized constraint is
      nonzero, so the random algebraic check catches the lie with constant
      probability.  Standard repetition reduces the remaining error.
\end{enumerate}
The essential role of the second prover is therefore not extra computation: it is
the enforcement of one globally consistent function behind locally requested
values.  This yields
\[
 \boxed{\MIP=\NEXP.}
\]

\subsection{Multilinearity test: the original BFL proof}

BFL work over the integer interval $I=\{0,\ldots,N-1\}$ and an oracle
$A:I^n\to\mathbb Q$, not with the later finite-field three-point test.  A line
is an axis-parallel copy of $I$.  It is \emph{correct} when the restriction of
$A$ is linear, and it is $\delta$-wrong when that restriction is not
$\delta$-close to any linear function.

\paragraph{The sampled-line test (Propositions 5.4--5.6).}
On a chosen line $\ell$, \textsc{Test}$_0$ samples $m_1+2$ points.  The first
two determine a linear function $h$; if $\ell$ is $\delta$-wrong, the remaining
samples all miss a disagreement with probability at most
$(1-\delta)^{m_1}<e^{-\delta m_1}$.  The full test samples $m_2$ lines in each
coordinate direction.  If more than an $\varepsilon$ fraction of the lines in
some direction are wrong, then the probability that no sampled line exposes
this is at most
\[
 e^{-\varepsilon m_2}+e^{-\delta m_1}.
\]
Thus, with $m_2=t/\varepsilon$ and $m_1=t/\delta$, acceptance lets the verifier
infer, with confidence at least $1-2e^{-t}$, that every direction contains
less than an $\varepsilon$ fraction of wrong lines.

\paragraph{Self-improvement (Lemma 5.10).}
Assume the wrong-line bound and suppose a multilinear $g$ already
$\Delta$-approximates $A$, where $\Delta\leq1/2$.  BFL partition $I^n$ into
bad points $B=\{A\neq g\}$, points on wrong lines, misplaced points on
approximately linear lines, and exceptional intersections where two distinct
linear functions agree.  For
$S=B\setminus(W\cup M)$, the expansion lemma gives
\[
 \left(1+\frac1{2n}\right)\mu(S)
 \leq \mu(\overline S)
 < \mu(S)+n(\varepsilon+\delta)+\frac nN.
\]
Consequently the same $g$ is actually $\varepsilon'$-close to $A$, with
\[
 \varepsilon'=3n^2\!\left(\varepsilon+\delta+\frac1N\right).
\]
This is why approximation error does not deteriorate at every inductive level.

\paragraph{Pasting hyperplanes (Lemma 5.11).}
Write a point as $(x,y)\in I\times I^{n-1}$.  Approximate linearity on most
$x$-lines produces functions $f_1,f_2$ for which
$xf_1(y)+f_2(y)$ is close to $A(x,y)$.  If a fraction at least
$\varphi=1/(10n)$ of the hyperplanes $x=\xi$ are already
$\beta$-approximately multilinear, with $\beta=1/10$, choose two good values
$\xi_1\neq\xi_2$.  Their multilinear approximants $g_1,g_2$ determine
\[
 \widetilde f_1(y)=\frac{g_1(y)-g_2(y)}{\xi_1-\xi_2},
 \qquad
 \widetilde f_2(y)=\frac{\xi_2g_1(y)-\xi_1g_2(y)}{\xi_2-\xi_1}.
\]
Hence $x\widetilde f_1(y)+\widetilde f_2(y)$ pastes the good hyperplanes into
one multilinear approximant; Lemma 5.10 then improves its error to
$\varepsilon'$.

\paragraph{Induction on aligned subspaces (Lemmas 5.12--5.14).}
BFL organize all aligned subspaces in a depth-$(n-1)$ tree: leaves are lines,
the root is $I^n$, and children are hyperplanes.  Correct lines are red.  A
node is red when almost all leaves below it are red and at least a
$\varphi$ fraction of its children are red.  The tree-colouring lemma pushes
redness to the root.  Inductively, red children are $\beta$-approximately
multilinear, and the Pasting Lemma makes their parent so as well.

The result is BFL Theorem 5.13: if each direction has less than an
$\varepsilon$ fraction of wrong lines, then $A$ is
$\varepsilon'$-approximately multilinear, where
\[
 \varepsilon'=3n^2\!\left(\varepsilon+\delta+\frac1N\right),
 \quad
 N\geq40n^2,
 \quad
 \delta\leq\frac1{400n^2},
 \quad
 \varepsilon\leq\frac1{800n^3}.
\]
Theorem 4.6 follows by choosing the parameters so that
$\varepsilon'\leq n^{-c}$.  This four-step argument---sample lines,
self-improve, paste, and colour the subspace tree---is the proof given in the
original paper.

\begin{keybox}
$\MIP=\NEXP$ implies that every $\mathsf{EXPTIME}$ language has a multiprover
interactive proof because $\mathsf{EXPTIME}\subseteq\NEXP$.  It does not prove
$\mathsf{EXPTIME}\not\subseteq\PSPACE$: whether
$\PSPACE=\mathsf{EXPTIME}$ is open.  Also, BFL is a classical theorem; allowing
entangled provers defines the very different class $\mathsf{MIP}^*$.
\end{keybox}

\section{GMW: zero knowledge}

For a verifier $V^*$, its view contains its input and auxiliary input, random
tape, and received messages.  A protocol is zero knowledge if for every
probabilistic polynomial-time $V^*$ there is an expected polynomial-time
simulator $S$ that uses no witness and produces
\[
 S(x)\ \begin{cases}
  =\operatorname{View}_{V^*}^P(x) &\text{perfect ZK},\\
  \approx_s\operatorname{View}_{V^*}^P(x)&\text{statistical ZK},\\
  \approx_c\operatorname{View}_{V^*}^P(x)&\text{computational ZK}.
 \end{cases}
\]
The definition quantifies over malicious verifiers, not only the honest one.
Simulation says that anything feasibly learned from the interaction could have
been generated from the statement alone.

\subsection{Warm-up: graph isomorphism}

For isomorphic $G_0,G_1$, the prover sends a random relabeling $H$ of one graph.
The verifier sends a challenge $b\in\{0,1\}$, and the prover opens an isomorphism
$G_b\cong H$.  A cheating prover able to answer both challenges gives an
isomorphism $G_0\cong G_1$.  A simulator guesses $b$, constructs $H$ from $G_b$,
runs $V^*$, and rewinds unless its challenge equals the guess.  Each attempt
succeeds with probability $1/2$, and conditioned on success the transcript has
the correct distribution.  Rewinding is the central simulation technique.

\subsection{The GMW protocol for 3-colouring}

Let $c:V\to\{1,2,3\}$ be a proper colouring and let the commitment scheme be
hiding and binding.
\begin{enumerate}
\item The prover chooses a random permutation $\pi$ of the three colours and
commits separately to every value $\pi(c(v))$.
\item The verifier chooses a uniformly random edge $uv$.
\item The prover opens only the two endpoint commitments; the verifier checks
that the openings are valid and the colours differ.
\end{enumerate}
Perfect completeness is immediate.  If the graph is not 3-colourable, every
fixed committed assignment has a monochromatic edge; binding prevents changing
it after the challenge.  A round therefore rejects with probability at least
$1/|E|$, and sequential repetition gives constant or negligible soundness error.

For a malicious verifier $V^*$ with arbitrary auxiliary input, the simulator
uses the following rewinding experiment.
\begin{enumerate}
\item Guess an edge $e=uv$ uniformly from $E$ and choose an ordered pair of
distinct colours uniformly from the six possibilities.
\item Commit to those colours at $u,v$ and commit to arbitrary dummy colours at
all other vertices; run $V^*$ on these commitments with a fixed sampled random
tape.
\item If $V^*$ asks for an edge different from $e$, discard the attempt and
restart.  Otherwise open the two commitments and output the complete view.
\end{enumerate}
For the honest verifier, the guessed edge matches with probability exactly
$1/|E|$.  For a malicious $V^*$, its challenge may depend on the commitment
string, so exact uniformity must not be assumed.  Instead, computational hiding
makes the challenge distribution computationally insensitive to which edge was
prepared.  The success probability is therefore inverse-polynomial, up to a
negligible term, and rewinding has expected polynomial running time (with the
standard cutoff if one wants a strict polynomial-time simulator).  Conditioned
on success, the opened ordered pair is uniform over the six distinct colour
pairs, exactly as for a uniformly permuted proper colouring.  The unopened
commitments need not encode a global colouring; a hybrid argument using
computational hiding replaces their dummy messages one at a time by the real
committed colours.  Thus the simulated and real views are computationally
indistinguishable.  Binding, separately, prevents a cheating prover from
adapting endpoint colours after seeing the edge.
Via reductions to 3-colouring, GMW obtains zero-knowledge proofs for all of
$\NP$ from the stated commitment/encryption assumption.

\textbf{Do not conflate properties.}  Binding supports soundness; hiding supports
zero knowledge.  Zero knowledge does not mean that the verifier learns nothing
at all---it learns that the statement is true.  Sequential composition is the
safe default under standard hypotheses; arbitrary parallel repetition need not
preserve zero knowledge.

\section*{Oral-exam checklist}
\begin{enumerate}
\item State completeness and soundness with the correct quantifier order.
\item Prove sum-check soundness by identifying the \emph{first} false polynomial.
\item For $\IP\subseteq\PSPACE$, discuss counters, recursion depth, and rational bit length.
\item Prove TQBF hardness with the universal selector that prevents formula duplication.
\item In $\PSPACE\subseteq\IP$, identify degree explosion and explain exactly how it is controlled.
\item Contrast Shamir's simple QBF with Shen's explicit $R_i$ operators.
\item Explain why noncommunication, not merely two answers, empowers $\MIP$.
\item Separate proof/argument, soundness/ZK, classical $\MIP$/$\mathsf{MIP}^*$, and
unrelativized/relativized claims.
\end{enumerate}

\section*{Primary references}
\begin{enumerate}
\item C.~Lund, L.~Fortnow, H.~Karloff, and N.~Nisan,
Algebraic Methods for Interactive Proof Systems, \emph{J.~ACM} 39 (1992), 859--868.
\item A.~Shamir, IP = PSPACE, \emph{J.~ACM} 39 (1992), 869--877.
\item A.~Shen, IP = PSPACE: Simplified Proof, \emph{J.~ACM} 39 (1992), 878--880.
\item L.~Babai, L.~Fortnow, and C.~Lund,
Non-Deterministic Exponential Time Has Two-Prover Interactive Protocols,
\emph{Computational Complexity} 1 (1991), 3--40.
\item O.~Goldreich, S.~Micali, and A.~Wigderson,
Proofs that Yield Nothing but Their Validity or All Languages in NP Have
Zero-Knowledge Proof Systems, \emph{J.~ACM} 38 (1991), 691--729.
\end{enumerate}

\end{document}